\chapter{結論}

本章では，本研究のまとめと今後の課題について述べる．

\section{まとめ}

本研究では，音楽嗜好の共有を通じた段階的な自己開示が受容性に与える影響について検討した．若者の深い自己開示を避ける傾向に対し，音楽嗜好の共有を通じた自己開示を提案し，カラオケ環境を想定した音楽共有セッションにおいて，段階的な自己開示を促す楽曲推薦システムを開発し，その効果を実証的に検証した．

実験からは，低レベルから高レベルの自己開示へと段階的に行うことで，高いレベルの自己開示に対する受容度が維持されることが確認された．システムによる段階的な推薦を受けた群は，後半フェーズで高いレベルの自己開示を促されているにも関わらず，推薦なし群と同等の高い受容度を維持していた．また，相手の音楽への認知度や興味度が低い場合でも，自己開示の受容度は大きく低下しないことが示された．これは，音楽を通じたコミュニケーションが，楽曲そのものへの興味の有無に関わらず，相手の経験や感情を受容できることを示唆している．さらに，音楽嗜好の共有は単なる趣味の共有を超えて，個人的な経験や価値観の自然な開示を促進することが明らかになった．

ヒアリングの結果からは，音楽をきっかけとしたコミュニケーションが長期的な友人関係においても新たな相互理解を促進されることが分かった．これらの知見は，音楽嗜好の共有が自己開示を促進し，人間関係の親密化に寄与する効果的な手段となり得ることを実証的に示している．


\section{今後の課題}

本研究には，いくつかの課題が残されている．歌唱を伴う実際のカラオケ環境での検証は重要な課題である．本実験では歌唱を除外した環境で基礎的な知見を得たが，実際のカラオケでは歌唱技術やプレッシャーが自己開示の深さに影響を与える可能性がある．今後は歌唱を含めた実験を行い，より実践的な知見を得る必要がある．

自己開示レベルの分類基準の精緻化も求められる．実験結果から，特に「思い入れ」と「自己開示レベル」の関係性について，より詳細な基準設定が必要であることが明らかになった．楽曲に対する個人的な意味づけをより適切に評価できる指標の検討も求められる．

また，本研究では20名という限られたサンプルサイズでの検証に留まっている．特に推薦システムの効果検証において，統計的な有意差を得るためには，より大規模なデータでの分析が必要である．様々な年齢層や関係性を持つペアでの検証を通じて，本手法の一般性を確認することも重要な課題である．

今後は，これらの課題に取り組むことで，音楽嗜好の共有を通じた自己開示の促進と受容に関する理解がさらに深まることが期待される．